% ===== 原文 pp.47–50（PDF p.57–60） =====

% [原文 p.47]
\para{约化后的作用量}

为了用重新标记后的场写出 Vafa-Witten 作用量，使用以下恒等式是有帮助的。设 $\chi_{\mu\nu}$ 与 $F_{\mu\nu}$ 是 4d 中的 2-形式，并令
\begin{gather*}
    \nu_r = \chi_{r4}, \qquad \tilde\nu_r = (\ast\chi)_{r4} \\
    \omega_r = F_{r4}, \qquad \tilde\omega_r = (\ast F)_{r4}
\end{gather*}
为把 $F$、$\ast F$、$\chi$ 与 $\ast\chi$ 限制到 $1,2,3$ 方向所得的 3d 1-形式，其中 $\ast$ 表示 4d Hodge 对偶。我们有
\[
    \sqrt{g}\left(\half\chi_{\mu\nu}F^{\mu\nu}\right)
    = \sqrt{g^{44}}\sqrt{g_W}\left(\nu_r\omega^r + \tilde\nu_r\tilde\omega^r\right)
\]
其中 $g_{\mu\nu}$ 是行列式为 $g$ 的 4d 度规，$(g_W)_{rs}$ 是行列式为 $g_W$ 的 3d 度规。类似地，我们有
\[
    F_{rs} = \sqrt{g^{44}}\sqrt{g_W}\epsilon_{rst}\tilde\omega^t
    = \sqrt{g^{44}}(\star\tilde\omega)_{rs}
\]
其中 $\star$ 是 3d Hodge 星。最后，设 $B_{\mu\nu}$ 是 4d 中的自对偶 2-形式，并令 $C_{\mu\nu}$ 为出现在 Vafa-Witten 局域化方程中的如下 2-形式：
\[
    C_{\mu\nu} = \half g^{\rho\tau}[B_{\mu\rho},B_{\nu\tau}] .
\]
人们发现
\[
    C_{r4} = (\ast C)_{r4} = \half\sqrt{g^{44}}\sqrt{g_W}\epsilon_{rst}[b^s,b^t]
    = \half\sqrt{g^{44}}(\star[b,b])_r
\]
其中 $b_r=B_{r4}$。对我们所取的圆周长，$g_{44}=g^{44}=1$。

% [Tutor Note: p.47 的第一个恒等式可用分块度规核对：$g_{r4}=0$，故 $\chi_{\mu\nu}F^{\mu\nu}=\chi_{rs}F^{rs}+2\chi_{r4}F^{r4}$，而 $F^{r4}=g^{44}g^{rs}F_{s4}=g^{44}F_{r4}$；把 $\chi_{rs}$ 用 3d Hodge 星写开即得 $\nu_r\omega^r$ 与 $\tilde\nu_r\tilde\omega^r$ 两项。取 $g_{44}=1$（本章约定）时 $\sqrt{g^{44}}=1$，可直接逐项比对。]

Vafa-Witten 作用量 (3.9) 用重新标记后的场写出后成为
\[
    S = \int_M \dd^4x\sqrt{g}\left(\cL_1+\cL_2+\cL_3+\cL_4\right)
    - \frac{\ii\tau}{4\pi}\int_M \Tr F\wedge F
\]
其中
\begin{align}
    \cL_1 &= \frac{1}{e^2}\delta_Q\Tr\left\{4\tilde\lambda^r\left(H_r - D_rA_4 + \partial_4A_r - (\star F)_r\right) + 2\bar\sigma\, \dd_A^\star\lambda\right\} \notag\\
    \cL_2 &= \frac{1}{e^2}\delta_Q\Tr\left\{\chi^r\left(\tilde P_r - 2D_r\phi + 2(\star \dd_Ab)_r\right) + \beta\left(P - 2\dd_A^\star b\right)\right\} \notag\\
    \cL_3 &= \frac{1}{e^2}\delta_Q\Tr\left\{4\tilde\lambda^r\left(-[b_r,\phi] + \half\star[b,b]_r\right) + 2\chi^rD_4b_r - 2\beta D_4\phi\right\} \notag\\
    \cL_4 &= \frac{1}{e^2}\delta_Q\Tr\left\{2\bar\sigma\left(-[\psi_r,b^r] - [\eta,\phi] - \half[\sigma,\tilde\rho] + D_4\rho\right)\right\} \tag{5.2}
\end{align}
其中 $\dd_A^\star\lambda = D_r\lambda^r$。若只是在上述作用量中保留常数模，所得 3d 作用量会含有形如
\[
    [A_4,\phi]
\]
的对易子项，而（正如我们一直强调的）由于 $A_4$ 的显式出现，这种项是定义不良的。

% [原文 p.48]
为了消去 $A_4$ 的出现，我们利用在写出给定上同调 TFT 的作用量时所拥有的广泛选择自由。也就是说，只要玻色动能项保持定义良好、且不改变理论的场含量与变分，我们就可以任意连续地形变作用量的 BRST-正合部分而不影响拓扑关联函数的值。我们写出一个新作用量 $\hat S$，把 $S$ 的 BRST-正合部分中出现的对易子 $[A_4,\phi]$ 换成
\[
    (2\pi)^{-1}\left(\operatorname{Ad}_{h^{-1}}\phi - \phi\right)
    = -(2\pi)^{-1}h^{-1}[h,\phi] .
\]
（后者在 $A_4\approx0$ 的邻域内化为前者。）可以用一族作用量
\[
    S_\alpha = (1-\alpha)S + \alpha\hat S
\]
在 $S$ 与 $\hat S$ 之间光滑插值，其中 $0\le\alpha\le1$，使得当我们把参数 $\alpha$ 由零调到一时，作用量被一个 BRST-正合的量光滑地形变。（非零 $\alpha$ 的效应是“开启”涉及 $A_4$ 对易子的高阶 BRST-正合项。）在 $\alpha=1$ 时，剩下的是等价的作用量 $\hat S$；应当强调，由于 $h$ 的出现，它并不是洛伦兹不变的。与 $S$ 不同，$\hat S$ 可以通过要求场不依赖 $x^4$ 坐标而安全地约化到 3d，于是得到
\[
    \hat S = \int_W \dd^3x\sqrt{g_W}\left(\hat\cL_1 + \hat\cL_2 + \hat\cL_3 + \hat\cL_4\right) + S_{top}
\]
其中
\begin{align}
    \hat\cL_1 &= \frac{1}{e^2}\delta_Q\Tr\left\{4\tilde\lambda^r\left(H_r + (2\pi)^{-1}h^{-1}D_rh - \star F_r\right) + 2\bar\sigma\, \dd_A^\star\lambda\right\} \notag\\
    \hat\cL_2 &= \frac{1}{e^2}\delta_Q\Tr\left\{\chi^r\left(\tilde P_r - 2D_r\phi + 2(\star \dd_Ab)_r\right) + \beta\left(P - 2\dd_A^\star b\right)\right\} \notag\\
    \hat\cL_3 &= \frac{1}{e^2}\delta_Q\Tr\left\{4\tilde\lambda^r\left(-[b_r,\phi] + \star[b,b]_r\right) - 2(2\pi)^{-1}\chi^rh^{-1}[h,b_r] + 2(2\pi)^{-1}\beta\, h^{-1}[h,\phi]\right\} \notag\\
    \hat\cL_4 &= \frac{1}{e^2}\delta_Q\Tr\left\{-2\bar\sigma\left([\lambda_r,b^r] + [\eta,\phi] - \half[\tilde\rho,\sigma] + (2\pi)^{-1}h^{-1}[h,\hat\rho]\right)\right\} \tag{5.3}
\end{align}
且 $S_{top}$ 是度量 4d 作用量中 4d 瞬子数的、依赖于耦合的项。

% [Tutor Note: p.48 的替换为何“在 $A_4\approx0$ 附近还原为 $[A_4,\phi]$”可直接展开检验：$h=\ee^{-2\pi A_4}$ 故 $h^{-1}=\ee^{2\pi A_4}\approx1+2\pi A_4$，于是 $\operatorname{Ad}_{h^{-1}}\phi=h^{-1}\phi h\approx\phi+2\pi[A_4,\phi]$，即 $(2\pi)^{-1}(\operatorname{Ad}_{h^{-1}}\phi-\phi)\approx[A_4,\phi]$；而 $(2\pi)^{-1}(\operatorname{Ad}_{h^{-1}}\phi-\phi)=-(2\pi)^{-1}h^{-1}[h,\phi]$（因为 $h^{-1}[h,\phi]=\phi-h^{-1}\phi h$），故所替换的两种写法彼此相等。符号是最易错处，务必按此展开核对。]
% [Tutor Note: p.48 的形变合法性依赖上同调 TFT 的基本性质：作用量的 BRST-正合部分可以任意形变而不改变 $\delta_Q$-闭可观测量的关联函数（见第 3 章 §3.1）。注意作者强调 $\hat S$ 不洛伦兹不变——这与“$h$ 的动能项是群流形上的 $\sigma$-模型项”一致，因为 $h$ 是群取值而非李代数取值。]

作用 $\delta_Q$ 后，人们发现玻色场 $h$ 的动能项取群流形 $\sigma$-模型映射的标准形式：
\[
    \Tr\{h^{-1}(D_rh)\, h^{-1}(D^rh)\}
\]
其中 $G$ 上的度规取为 $\Tr$ 在单位元处诱导的左不变度规。

% [原文 p.49]
把辅助场积掉后，人们发现理论局域化在 BPS 方程
\begin{align}
    F - (2\pi)^{-1}h^{-1}\star \dd_Ah + [\star b,\phi] - [b,b] &= 0 \notag\\
    \star \dd_Ab - \dd_A\phi - (2\pi)^{-1}h^{-1}[h,b] &= 0 \notag\\
    \dd_A^\star b - (2\pi)^{-1}h^{-1}[h,\phi] &= 0 \tag{5.4}
\end{align}
上，以及涉及 $\sigma$ 的方程上。

上述场含量、作用量与变分正是靶空间为 $\fg$ 的带规范场的 3d A-模型的。规范扇区并不精确是 B.2.1 节所述的 3d A-型规范理论，而是其中李代数取值的标量场 $\varsigma$ 被群取值的场 $h$ 取代的一个修改。

\para{瞬子项}

我们希望把上述作用量中的瞬子项按 3d 的语言来解释。分析分两种情形：$G=U(1)$ 与 $G$ 为单的情形。

在 $G=U(1)$ 的情形，我们把瞬子项的约化写成
\begin{equation}
    S_{top} = -\frac{\ii\tau}{2}\int_W F\wedge \dd A_4 \tag{5.5}
\end{equation}
注意，正是 $A_4$ 的周期性使这个拓扑项一般而言非平凡。

% [Tutor Note: p.49 式 (5.5) 的非平凡性来自 $A_4$ 是圆取值标量：若 $A_4$ 是普通实标量，则 $\dd A_4$ 恰当，$\int_W F\wedge \dd A_4=\int_W\dd(F A_4)=0$（$W$ 无边）；而当 $A_4$ 周期为 $1$ 时 $\dd A_4$ 可有非零整数周期，该积分一般不为零——这正是 4d 拓扑项在 3d 中的残余。取 $A_4\to A_4+\ii m$（$m\in\Z$）时 $\dd A_4$ 不变，可作为该表达式周期性的初步自洽检验。]
% [Tutor Note: p.49 式 (5.4) 的第一式是 2-形式方程、第二式 1-形式、第三式 0-形式，可由每一项的指标数核对；令 $h\to1$（即 $A_4\to0$）时含 $h^{-1}[h,\cdot]$、$h^{-1}D_rh$ 的项消失，(5.4) 退化为通常的 Vafa–Witten 局域化方程。这是判断 $2\pi$ 归一化与各项次数是否抄错的最快方法。]

在 $G$ 为单的情形，没有这种形变项可用来度量 $W$ 上主丛的拓扑类。事实上，由于我们已把 $W$ 取为紧致的，该项自动为零。更一般地，我们希望考虑插入无序算子并允许 $W$ 有边界的可能性。例如，在第 6 章中我们将把 4d 理论的面算子的 2-范畴描述为 3d 边界条件的 2-范畴。然而，3d 形变的缺乏表明，凡是可以纯粹用 3d 语言理解的理论性质都将与 $\tau$ 无关，例如面算子的 2-范畴结构。正如我们在第 4 章提到的，这个 2-范畴暗中具有辫子幺半结构（对应于面算子的碰撞），而这一额外结构可能依赖于 $\tau$。但由于我们将不考察辫子幺半结构，我们从 3d 紧化理论的作用量中丢弃瞬子项。

% [原文 p.50]
\section{GL-扭转理论在 \texorpdfstring{$t=0$}{t=0} 时的 \texorpdfstring{$S^1$}{S1} 紧化}

现在转向 GL-扭转理论的约化，先处理 $t=0$ 的情形。对每个 $t$ 值，分析都将类似地进行。Vafa-Witten 理论与 $t=0$ 理论最后在三维中约化为同一个理论。

Vafa-Witten 扭转理论与 $t=0$ 在流形 $W\times S^1$ 上相同，这一事实可以如下理解。回忆 2.4.1 节：对 $\cN=4$、$d=4$ 的 SYM 作 Vafa-Witten 扭转后，标量超荷 $Q$ 取自左手 $Q$ 之中。另一方面，作 GL-扭转也恰在 $t=0$ 时给出一个左手标量超荷 $Q_l$；确实，它可以取为同一个。于是 Vafa-Witten 与 $t=0$ 两种理论仅仅在场于 4d 旋转群下的变换方式上不同；例如，Vafa-Witten 理论有一个自对偶 2-形式 $B_{\mu\nu}$ 与 0-形式 $\phi$，而 $t=0$ 理论有一个 1-形式 $\phi_\mu$。然而，这些场在 3d 旋转子群下的变换是相同的；例如，两个 3d 约化理论都有一个 3d 1-形式 $b_r$ 与一个 3d 0-形式 $\phi$。

明确地，回忆 2.5.2 节中 $t=0$ 时 GL-扭转理论的作用量：
\[
    S = \int_M \dd^4x\sqrt{g}\left(\cL_1+\cL_2\right)
    - \frac{\bar\tau}{4\pi\ii}\int_M \Tr F\wedge F
\]
其中
\begin{align}
    \cL_1 &= \frac{1}{e^2}\delta_Q\Tr\left\{\chi_{\mu\nu}H^{\mu\nu}
        - 2\chi_{\mu\nu}^+\left(F-\half[\phi,\phi]\right)^{\mu\nu}
        + 2\chi_{\mu\nu}^-\left(\dd_A\phi\right)^{\mu\nu}\right\} \notag\\
    \cL_2 &= \frac{1}{e^2}\delta_Q\Tr\left\{-\tilde\eta\left(P - 2\dd_A^\star\phi\right)
        + \bar\sigma\left(2\dd_A^\star\psi + 2[\tilde\psi_\mu,\phi^\mu] + [\sigma,\eta]\right)\right\} . \tag{5.6}
\end{align}

% [Tutor Note: p.50 式 (5.6) 的拓扑项系数 $-\bar\tau/(4\pi\ii)=+\ii\bar\tau/(4\pi)$，与 §5.2 的 $-\ii\tau/(4\pi)$ 在 $\theta=0$ 时给出同样的 $1/e^2$ 项，但 $\theta$ 项符号相反（等价于 $\tau\to-\bar\tau$）。比较 (3.9) 与 (5.6) 时不要误判为笔误：这反映第 2、3 章对 $\theta$ 的符号或以 $F$ 的定向约定不同。]
% [Tutor Note: p.50 关于“紧化后同一理论”的论证要点是：两种扭转给出同一个左手标量超荷 $Q_l$，因此 BRST 复形相同；差别只在 4d 旋转群表示（$B_{\mu\nu}\oplus\phi$ 对 $\phi_\mu$），而二者的 3d 旋转子群表示一致（$b_r\oplus\phi$）。核对方法：把 (5.6) 的场按 $r,s,t=1,2,3$ 与 $4$ 分解，与 §5.2 的 (5.1) 逐项配对，即得下面两张表。]
% [Tutor Note: p.50 原文 "under the the 3d rotation subgroup" 重复了定冠词 "the"，属笔误，应为 "under the 3d rotation subgroup"。]

我们按 3d 语言重新标记 GL-扭转理论的场如下（上标 $[0]$ 表示 $t=0$ 时 GL-扭转理论中相应的场，如我们在第 2 章中所写）。规范扇区的场为
\begin{gather*}
    A_r = A_r^{[0]}, \qquad \lambda_r = \ii\psi_r^{[0]} \\
    \varsigma = A_4^{[0]}, \qquad \rho = \ii\psi_4^{[0]} \\
    \sigma = -\ii\sigma^{[0]}, \\
    \bar\sigma = -\ii\bar\sigma^{[0]}, \qquad \tilde\rho = \eta^{[0]} \\
    \tilde\lambda_r = \chi_{r4}^{+[0]}, \qquad H_r = H_{r4}^{+[0]}
\end{gather*}
物质扇区的场为
\begin{gather*}
    \phi = -\phi_4^{[0]}, \qquad \eta = \ii\tilde\psi_4^{[0]}, \\
    b_r = -\phi_r^{[0]}, \qquad \psi_r = \ii\tilde\psi_r^{[0]} \\
    \beta = \tilde\eta^{[0]}, \qquad P = -P^{[0]} \\
    \chi_r = 2\chi_{r4}^{-[0]}, \qquad \tilde P_r = 2H_{r4}^{-[0]} .
\end{gather*}
容易验证，重新标记后的场的 BRST 变分由上面的 (5.1) 给出。
